% SPDX-License-Identifier: AGPL-3.0-or-later
%
% CliffordClock methods-and-validation paper (WP15 draft).
%
% DRAFT ONLY. The project owner is the sole author of record, reviews
% every claim below, and decides whether/when/where to submit this
% document. No agent has submitted, posted, or otherwise distributed
% this paper anywhere. The author block was filled in by the owner;
% an acknowledgments section can be added by the owner at finalization.
%
% Build: see paper/Makefile (agent build/style notes are kept in the
% project's internal build-notes record, untracked). Every numeric value below
% that originates from the codebase is \input from paper/generated/*.tex,
% produced by paper/figures/*.py against the real cliffordclock pipeline
% and benchmarks/ code; never hand-typed.

\documentclass[aps,pra,superscriptaddress,groupedaddress,nofootinbib]{revtex4-2}

\usepackage{amsmath}
\usepackage{amssymb}
\usepackage{graphicx}
\usepackage{booktabs}
\usepackage{xcolor}
\usepackage{hyperref}
\usepackage{placeins}
\usepackage{float}

% Loosen LaTeX's default float-placement fractions: this paper carries
% more floats (5 figures, 3 tables) than typical default parameters
% budget for well, and the defaults were causing large deferred blank
% regions rather than tightly packed pages.
\renewcommand{\topfraction}{0.9}
\renewcommand{\bottomfraction}{0.85}
\renewcommand{\textfraction}{0.07}
\renewcommand{\floatpagefraction}{0.75}
\setcounter{topnumber}{4}
\setcounter{bottomnumber}{4}
\setcounter{totalnumber}{8}

% --- Generated numeric macros (paper/figures/*.py; regenerate via `make figures`). ---
\input{generated/npl_values.tex}
\input{generated/step_size_values.tex}
\input{generated/showcase_values.tex}
\input{generated/bbr_jila_values.tex}
\input{generated/bbr_temperature_values.tex}
\input{generated/bothwell_values.tex}
\input{generated/roos_values.tex}
\input{generated/budget_slice_values.tex}

\usepackage{orcidlink}

\begin{document}

\title{CliffordClock: a Cl(1,3) spacetime-algebra pipeline for field-gradient
and DC-Stark frequency-shift budgeting in optical lattice clocks}

\author{Michael Rudolph\,\orcidlink{0000-0002-0495-5723}}
\affiliation{Velar Scientific}
\email{michael.rudolph@velarscientific.com}

\date{\today}

\begin{abstract}
Optical lattice and ion clocks target systematic-uncertainty budgets at
the $10^{-18}$ level, where DC-Stark shifts from residual stray electric
fields, and the spatial variation of those fields across an extended
atomic cloud, are a routinely underdocumented systematic. We present CliffordClock, an
open-source pipeline that takes an imported (CAD/FEA-exported) or
synthetic electric-field model with real spatial structure, propagates it
through an atomic ensemble (a classical Monte Carlo cloud moving through
the field, or a Hermite--Gauss lattice quadrature), and reports the mean
fractional-frequency shift together with the full inhomogeneous-shift
budget a spatially varying field produces: the shift's ensemble spread,
the resulting dephasing time $T_2^*$, and the clock transition's spectral
line profile. A field-gradient map alone hands a lab none of these; the
pipeline computes them end to end, from field export to a reported,
uncertainty-quantified budget. The physical coupling is an
ordinary scalar second-order DC-Stark pivot,
$\Delta\nu/\nu_0 = -(\Delta\alpha/2)|E(\mathbf{r})|^2/(h\nu_0)$, using
literature differential polarizabilities; this scalar evaluation fully
handles spatially varying fields on its own. Three evaluation modes share
this physics: a fast analytic mode for lattice-clock motional states
(exact, no time integration), a trajectory mode that propagates a
Monte Carlo ensemble's actual motion through a spatially varying field,
and a general Cl(1,3) spacetime-algebra rotor mode, verified to agree
with the scalar modes to floating-point precision on every case this
paper models and serving as the formalism this project is built on for
physics the scalar modes cannot express
(Sec.~\ref{sec:limitations}'s roadmap). We
validate against four closed-form self-consistency cases, four literature
known-answer cases ($^{87}$Sr, $^{171}$Yb DC-Stark and second-order-Doppler
formulas, $\mathrm{rtol}=10^{-10}$ or better), a direct rotor/scalar
cross-check, and two zero-free-parameter \emph{reproducibility}
cases: NPL's published DC-Stark shift band, and Bothwell et al.'s
published mm-scale redshift slope. A blind prediction, against data the
pipeline's inputs did not already combine, would be the stronger claim;
the project currently reports zero blind-prediction cases, and
producing one is the concrete next validation goal.
A showcase scenario propagates a Monte Carlo Sr-87 cloud through a
finite-difference-solved chamber field with genuine curvature (two
asymmetric electrodes plus a patch-potential wall spot), reporting a mean
fractional shift of $\ShowcaseMeanShift$ with a $\ShowcaseShiftSpread$
ensemble spread and $T_2^*=\ShowcaseTtwoStarUs\,\mu$s, computed
identically (to $\ShowcaseMeanShiftRelDiff$ relative) through both the
trajectory and rotor evaluation modes. Code, tests, and
figure-generation scripts are public under the AGPLv3.
\end{abstract}

\maketitle

\section{Introduction}
\label{sec:intro}

Optical atomic clocks based on neutral-atom optical lattices and trapped
ions have reached fractional-frequency systematic uncertainties at or
below $10^{-18}$~\cite{Ludlow2015,Aeppli2024,Jia2026}. At this level, a
clock's DC-Stark shift budget, the frequency shift from a
transition's differential static scalar polarizability $\Delta\alpha$
coupling to the square of a residual stray electric field at the atoms,
is no longer a formality: it is routinely one of several systematics
that must be independently bounded, and its spatial \emph{gradient}
across an extended atomic ensemble (a Ramsey-Bord\'e cloud, a
multi-site optical lattice, or an ion crystal) can broaden the clock
transition line and add to the shift's own uncertainty even after the
mean field is nulled. Groups building or upgrading a clock apparatus
typically characterize this budget with one of three approaches: a
general-purpose finite-element solver (COMSOL, SIMION, or similar) whose
output is manually carried into a hand or spreadsheet Stark-shift
calculation; a lab's own unpublished script, rewritten by each new
graduate student; or, for detailed treatments of specific atomic
properties (Rydberg-state Stark maps, polarizability calculations), a
dedicated tool such as ARC~\cite{Sibalic2017}, in the broader spirit
of validated, community-maintained physics simulation packages such as
QuTiP~\cite{Lambert2024} for open quantum systems. None of these directly
close the loop from an imported field-gradient map (an FEA export, or a
hand-measured stray-field profile) to a reported, uncertainty-quantified
fractional-frequency shift and dephasing time for a specific trapped
ensemble and interrogation protocol. That is the workflow gap this
project addresses.

CliffordClock is a Python pipeline that takes an electric-field model
(a closed-form synthetic field, or an imported CSV/COMSOL grid with the
real spatial structure a CAD/FEA export or a hand-measured survey
actually has) and an atomic ensemble specification (species, trap,
motional regime, interrogation time), and closes the loop that field map
alone leaves open: it reports the resulting fractional frequency shift,
its statistical uncertainty, \emph{and} the full inhomogeneous-shift
budget the field's spatial variation produces across the ensemble. That
budget is the shift's spread, the inhomogeneous dephasing time $T_2^*$,
and the clock transition's spectral line profile. A lab with only a field-gradient map
and a hand or spreadsheet calculation gets a single mean-shift number at
best; this pipeline turns the same map into the dispersion budget that
map implies, for a specific trapped ensemble and interrogation protocol,
without re-deriving the perturbation theory or re-implementing a
numerically stable field-squared evaluation from scratch. Sec.~\ref{sec:showcase}
demonstrates this end to end on a chamber-scale field with genuine
curvature. The target regime is the $10^{-18}$-level DC-Stark and
gradient budgeting relevant to current optical lattice clocks (worked
here for $^{87}$Sr and $^{171}$Yb, the two $J{=}0\!\to\!J{=}0$ species
with literature-cited differential polarizabilities in the project's
species registry).

Physically, this is an ordinary second-order (quadratic) Stark shift:
optical-clock states carry no permanent electric dipole moment, so the
field couples via the transition's differential static scalar
polarizability, and this scalar coupling
fully handles a spatially varying field on its own: nothing about a
field gradient or curvature requires anything beyond it. What this
project adds underneath is a second, independently re-derivable
formulation of the same physics in the Cl(1,3) spacetime algebra of
proper time and spin connections~\cite{Hestenes2015,DoranLasenby2003}: a
\emph{rotor} engine (a rotor is a rotating geometric object representing
the atom's internal clock phase; advancing it along the atom's
\emph{worldline}, its path through space and time, is the general
relativistic-kinematics calculation this project is built on) capable of
representing boost/rotation effects exactly, with no small-angle or
single-order perturbative truncation in its own exponential map.
Section~\ref{sec:model} describes the three evaluation modes and what
each is for. The scalar modes are performance tiers of one shared
physical model, and the rotor mode computes the same physics through
the Cl(1,3) algebra; the two agree to floating-point precision on every
case this paper computes. The rotor engine earns its place through the
roadmap physics of Table~\ref{tab:scope}. Magnetic couplings,
internal-state structure, and transport in accelerating frames all
carry directional or geometric content, and none of them can be written
as a simple shift value at each point in space the way the DC-Stark and
blackbody terms can; representing them takes the rotor's richer
mathematical object.

\section{Physical model}
\label{sec:model}

\subsection{The scalar DC-Stark pivot}
\label{sec:scalar-pivot}

The transition frequency at a point $\mathbf{r}$ in the trap, relative to
its unperturbed value $\nu_0$, is written as a dimensionless
\emph{pivot} $P(\mathbf{r})$, the local fractional rate factor
multiplying proper time, such that the local proper-time rate is
$d\tau = P(\mathbf{r})\sqrt{1-v^2/c^2}\,dt$. For the physical DC-Stark
coupling used throughout this paper's validated results,
\begin{equation}
  P(\mathbf{r}) - 1 \;=\; \frac{\Delta\nu(\mathbf{r})}{\nu_0}
  \;=\; -\frac{\Delta\alpha}{2}\,\frac{|E(\mathbf{r})|^2}{h\nu_0} ,
  \label{eq:e14b}
\end{equation}
with $\Delta\alpha = \alpha_{\rm excited}-\alpha_{\rm ground}$ the
transition's differential static scalar polarizability and $\nu_0$ the
clock transition frequency (code cross-reference: CONVENTIONS.md
E14b). Equation~\eqref{eq:e14b} is the same textbook second-order
Stark formula used throughout the optical-clock
literature~\cite{Middelmann2012,Sherman2012,Bowden2017}, and every
evaluation mode in this paper (Sec.~\ref{sec:model-modes}) computes it
\emph{directly} from the total field $E(\mathbf{r})$ at each point.
There is no baseline/perturbation split at this stage (that split,
$E_0$/$\delta E$, belongs to the field smoother upstream,
Eqs.~\eqref{eq:e11}--E12 below), and no restriction to a spatially
uniform field anywhere in this formula: a field with a real gradient or
curvature across the ensemble is evaluated at each atom's own position,
exactly, by the scalar pivot alone. Sec.~\ref{sec:methods} shows why
this direct evaluation is numerically safe and how that safety is
pinned by an adversarial regression test. The
instantaneous fractional rate perturbation including kinematic
(second-order Doppler) time dilation is
\begin{equation}
  \delta\tilde\omega(\mathbf{r},\mathbf{v}) \;=\;
  P(\mathbf{r})\sqrt{1-v^2/c^2} - 1 ,
  \label{eq:e21}
\end{equation}
(E21) integrated over an atom's trajectory to give its accumulated
fractional phase, and averaged over the ensemble (uniform weights for a
classical Monte Carlo sample; for a lattice motional-state ensemble,
Gauss--Hermite quadrature weights, a small set of weighted sample points
that exactly average polynomial functions over the atomic wavefunction)
to give the reported $\langle\Delta\nu/\nu_0\rangle$ (E22--E23).
For a $J{=}0\!\to\!J{=}0$
clock transition (the lattice species $^{87}$Sr and $^{171}$Yb, and the
$J{=}0\!\to\!J{=}0$ ion clocks $^{27}$Al$^+$ and $^{115}$In$^+$, all in
the species registry) the scalar $\Delta\alpha$ is the whole DC-Stark
story. An ion clock whose upper state has $J\ge1$ additionally carries
a permanent electric-quadrupole moment coupling to the field
\emph{gradient} at first order; that term is implemented and described
in Sec.~\ref{sec:quadrupole}, while the tensor-polarizability
correction to the second-order Stark shift itself remains roadmap
physics (Sec.~\ref{sec:limitations}). A second, linearized coupling,
$P(\mathbf{r})-1 = \delta E(\mathbf{r})\!\cdot\!\boldsymbol\mu / (m_ec^2)$
(E14a, an explicit effective dipole moment $\boldsymbol\mu$), is also
implemented; Eq.~\eqref{eq:e14b} provably reduces to it under a
bias-field linearization (CONVENTIONS.md E14b "Bridge identity"). E14a
is not part of this paper's physical claims about a real DC-Stark
shift; it is used only for the closed-form self-consistency cases of
Sec.~\ref{sec:validation} (Table~\ref{tab:validation}, rows V1--V4),
where an explicit, hand-tunable coupling constant is the useful property.

\subsection{Three evaluation modes, one physical model}
\label{sec:model-modes}
\label{sec:rotor-verification}

Eq.~\eqref{eq:e14b}'s pivot is evaluated by three user-facing evaluation
modes, all computing the same physical quantity, that trade generality for
speed in a controlled, verified way:

\begin{itemize}
\item \textbf{Fast analytic mode} (\texttt{integration.mode=fast\_path},
  the default for lattice-clock motional states): for a static
  Gauss--Hermite quadrature ensemble (Sec.~\ref{sec:scalar-pivot}) in a
  time-independent field, the
  phase accumulated over an interrogation time $T$ is exactly
  $\Delta\Phi_q = \delta\tilde\omega_q\,\tilde T$ per quadrature node
  (E29), an algebraic evaluation with no time integration whose cost is
  therefore independent of $T$.
\item \textbf{Trajectory mode} (\texttt{integration.mode=direct} for
  \texttt{ensemble.regime=classical}, the default there): a Monte Carlo
  ensemble's atoms are propagated on their actual classical trajectories
  through the field (velocity-Verlet dynamics in the trap, E21's pivot
  evaluated along the way), so a spatially varying field's real structure
  (gradients, curvature, any shape a CAD/FEA export carries) is
  sampled exactly where each atom's trajectory visits it. This is the
  mode Sec.~\ref{sec:showcase}'s showcase runs.
\item \textbf{Rotor mode} (\texttt{integration.mode=worldline}): the full
  Cl(1,3) spacetime-algebra engine below, advancing each atom's Lorentz
  frame along its worldline (its path through space and time).
\end{itemize}

\textbf{For every case this paper models, the scalar modes (fast analytic
and trajectory) and the rotor mode agree to floating-point precision.}
This is a directly verified statement: the scalar and
rotor accumulators are run on \emph{identical} trajectories/nodes and
compared head to head, both in the codebase's standing test suite
(\path{tests/test_integrator_stark_rotor.py}) and in this paper's own
showcase computation (Sec.~\ref{sec:showcase}, Fig.~\ref{fig:showcase}),
where the two evaluations of the same Monte Carlo ensemble agree to
$\ShowcaseMeanShiftRelDiff$ relative in the mean shift and
$\ShowcaseMaxPhaseDiff$ absolute in per-atom accumulated phase. The
scalar modes exist purely for speed and are not an approximation of a
"more correct" rotor calculation happening underneath: for every physical
coupling this project currently implements, the rotor result and the
scalar result are the same number, computed two independent ways.

The rotor engine is nonetheless the general formalism this project is
built on, and its distinct necessity is not hypothetical: it begins
exactly where a physical effect stops being expressible as a single
number ($\delta\tilde\omega$) per point in space, which is where every
item in Table~\ref{tab:scope}'s roadmap column, and Sec.~\ref{sec:limitations}'s
extension list (magnetic couplings, internal-state structure,
transport-induced geometric phases), lives. The trapped atom's local
Lorentz frame is advanced along its worldline by
\begin{equation}
  \frac{dR(\tilde\tau)}{d\tilde\tau} = -\tfrac12\,\Omega(\mathbf{r}(\tilde\tau))\,R(\tilde\tau) ,
  \label{eq:e17}
\end{equation}
(E17, times in Compton units $\tilde\tau=\tau/\tau_C$,
$\tau_C=\hbar/m_ec^2$; $R$ is a \emph{rotor}, a rotating geometric
object representing the atom's internal clock phase) with the
\emph{interaction bivector} $\Omega = \omega_{\rm boost} + I\,\omega_{\rm
rot}$ (E18; a bivector is the geometric-algebra object representing an
oriented plane, here the plane of the internal clock's rotation), built
from the pivot's spatial gradient (a boost/frame-tilt bivector,
$\omega_{0k}=\partial_k\ln P(\mathbf{r})$, E16) and the pivot's own
internal rotation rate. The rotor equation is integrated with an
exponential-midpoint stepper, exactly rotor-group-preserving at any step
size (E19), and the phase it extracts by composing successive
$\exp(\Omega)$ rotors is compared against the primary scalar phase of
Eq.~\eqref{eq:e21} as a standing cross-check (E24): agreement is required
at first order in the boost bivector, with a genuine second-order
($O(\omega_{\rm boost}^2)$) divergence between the two paths (visible
only for gradients far steeper than any realistic lattice-clock trap)
expected by construction, the rotor result being the more complete one
in that regime.

An $\Omega$-bivector construction for Eq.~\eqref{eq:e14b}
(\texttt{build\_omega\_stark}, instantiating E16/E18 for the physical
DC-Stark coupling) drives \texttt{integration.mode=worldline}'s true
Cl(1,3) rotor directly under \texttt{coupling.type=stark\_dc}. The
scalar/rotor equivalence under the physical coupling is verified
directly: the rotor-vs-scalar test suite
(\path{tests/test_integrator_stark_rotor.py}) checks first-order
agreement at realistic field/velocity scale, the permitted
$O(\omega_{\rm boost}^2)$ divergence under an exaggerated boost, a
uniform-field null, a $v=0$ static-node identity, a non-symmetric-field-
gradient orientation guard (ruling out an axis-transposition error in the
spin-connection contraction), and step-halving convergence of the rotor
accumulator under a genuinely time-varying interaction bivector along a
moving trajectory (measured $\mathcal{O}(d\tilde\tau^2)$, ratio
$\approx4$ per halving, against an endpoint-evaluated stand-in that
measures $\approx2$); the NPL reproducibility case of
Sec.~\ref{sec:validation} is independently re-run end to end through this
rotor path, reaching the identical \textit{MET} verdict and, to float64-relative
precision, the identical predicted band as the scalar case; and
Sec.~\ref{sec:showcase}'s showcase re-runs a genuinely moving Monte Carlo
ensemble through both paths at once.

\subsection{The blackbody-radiation pivot term}
\label{sec:bbr-pivot}

A second contribution enters the same scalar pivot $P(\mathbf{r})$ that
Eq.~\eqref{eq:e14b} defines for the DC-Stark shift. Any vacuum chamber
above absolute zero glows faintly in the infrared: its walls radiate
thermal photons, the way a warm object radiates heat across a
distance. That glow is itself an electric field, and it
couples to the atom's clock states through the very same differential
polarizability $\Delta\alpha(0)$ the DC-Stark term already uses. The
blackbody-radiation (BBR) pivot term is
\begin{equation}
  (P-1)_{\rm BBR} = \frac{\Delta\nu_{\rm stat}\,(T/T_0)^4 +
  \Delta\nu_{\rm dyn}(T)}{\nu_0} ,
  \label{eq:e32}
\end{equation}
referenced to a room-temperature baseline $T_0=300$~K (E32). The
static coefficient is $\Delta\nu_{\rm stat} =
-(\Delta\alpha(0)/2h)\langle E^2\rangle_{T_0}$, under the full
energy-density convention $\langle E^2\rangle_T = (4\sigma_{\rm
SB}/\epsilon_0 c)\,T^4$~\cite{Middelmann2012,Lisdat2021}. The equation
carries no separate leading minus sign: $\Delta\nu_{\rm stat}$ is
already negative, and that alone is what makes the shift negative, in
exactly the same $P-1=\Delta\nu/\nu_0$ convention
Eq.~\eqref{eq:e14b} established.

The DC-Stark and BBR contributions then combine the simplest way
possible: they add, $P(\mathbf{r})-1 = (P-1)_{\rm
stark}(\mathbf{r}) + (P-1)_{\rm BBR}$ (E33). One might worry about a
cross term between the static trap field and the fluctuating thermal
field, but it vanishes exactly: the thermal field is isotropic, so its
vector average $\langle E_{\rm BBR}\rangle$ is zero and the cross term
averages away with it. The one
product term that does survive, $(P-1)_{\rm stark}\cdot(P-1)_{\rm
BBR}\sim10^{-33}$, is far below anything this paper reports. The bath
itself needs no stochastic treatment either: a roughly one-second
interrogation averages over $\sim10^{14}$ independent femtosecond-scale
bath cycles, so the atom experiences the deterministic expectation
$\langle E^2\rangle$, the same deterministic treatment every published
BBR evaluation uses.

$\Delta\nu_{\rm dyn}(T)$ is not a truncated $1/T$ power series. Lisdat,
D\"orscher, Nosske, and Sterr proved this Taylor expansion has zero
convergence radius for Sr's dominant BBR-coupled
transition~\cite{Lisdat2021}, and the historical three-term truncation
left published Sr corrections low by a documented $4.7\times10^{-18}$.
The registry instead fits $\Delta\nu_{\rm dyn}(T)$ as a small polynomial
to the exact Planck integral. For $^{87}$Sr this uses the PTB group's
rescaling of the Lisdat shape onto Aeppli et al.'s more precise
anchor~\cite{Nosske2025,Aeppli2024}, cross-verified to $10^{-19}$ for
$T\le300$~K (Sec.~\ref{sec:jila-bbr}'s benchmark). Because the bath is
isotropic and every registered species is $J{=}0\!\to\!J{=}0$, the term
reaches the rotor engine only through this scalar pivot, with zero
spin-connection gradient for uniform $T$ (Sec.~\ref{sec:limitations}).

Figure~\ref{fig:bbr-temperature} shows this pivot across its full
validity window, with the published dataset it inherits its shape from
laid directly on top: the blue curve is Eq.~\eqref{eq:e32} evaluated by
the real engine from $\BbrTempWindowLo$ to $\BbrTempWindowHi$~K with a
coefficient-uncertainty band, the red points are the Lisdat,
D\"orscher, Nosske, and Sterr open dataset's own tabulated total BBR
shift~\cite{Lisdat2021} converted to the same fractional units via this
project's own registry clock frequency, and the black star marks the
operating point Aeppli et al.\ at JILA published independently
(Sec.~\ref{sec:jila-bbr}).
Because the registry's dynamic-term polynomial is the PTB-2025
rescaling of this same dataset's fit shape onto Aeppli et al.'s more
precise 2024 anchor, the curve and the dataset points trace the same
shape by construction, and the lower panel's residual is not
measurement noise: it is the real, expected offset that revision
introduces, growing from a few parts in $10^{20}$ near
$\BbrTempWindowLo$~K to $\BbrMaxAbsResidual\times10^{-19}$ near
$\BbrMaxAbsResidualT$~K, with a $\BbrOffsetAtThreeHundredK\times10^{-19}$
offset at the 300~K point itself, everywhere larger than the dataset's
own stated uncertainty. This figure visualizes the same
arithmetic-reproduction relationship Sec.~\ref{sec:jila-bbr} already
documents at Aeppli et al.'s single operating point, extended across the full
window; it introduces no new validation case and does not change this
paper's validation headline.

\begin{figure}[htbp]
\centering
\includegraphics[width=0.46\linewidth]{figures/fig5_bbr_temperature.pdf}
\caption{The Sr-87 BBR pivot (Eq.~\eqref{eq:e32}) across its
$\BbrTempWindowLo$-$\BbrTempWindowHi$~K validity window (blue curve,
shaded coefficient-uncertainty band), overlaid with the Lisdat et al.\
2021 open dataset's own tabulated total BBR shift (red points,
converted to fractional units via this project's own registry clock
frequency) and Aeppli et al.'s published operating point (black star,
Sec.~\ref{sec:jila-bbr}). Bottom panel: the engine-minus-dataset
residual, in units of $10^{-19}$, with error bars from the dataset's own
stated uncertainty. Because the registry's dynamic-term coefficients are
the PTB-2025 rescaling of this dataset's own fit shape onto a more
precise later anchor (Sec.~\ref{sec:bbr-pivot}), the residual's growing
structure is the expected, real signature of that anchor revision.
Regenerated by
\texttt{paper/figures/fig5\_bbr\_temperature.py} against the real
\texttt{cliffordclock.integrator.omega.bbr\_pivot\_perturbation} and the
committed dataset fixture
\texttt{benchmarks/data/lisdat2021\_dataset/}.}
\label{fig:bbr-temperature}
\end{figure}

\subsection{The gravitational-redshift pivot term and extended samples}
\label{sec:grav-pivot}

General relativity contributes the next pivot term, and it needs no
electromagnetic field at all: two clocks at different heights in
Earth's gravity tick at different rates. For a site or atom at height
$h(\mathbf{r}) = \hat u\cdot\mathbf{r}$ along a configured up direction
$\hat u$, relative to a configured reference height $h_{\rm ref}$, the
weak-field clock-rate factor enters the pivot as
\begin{equation}
  (P-1)_{\rm grav}(\mathbf{r}) \;=\;
  \frac{g\,\bigl(h(\mathbf{r}) - h_{\rm ref}\bigr)}{c^2} ,
  \label{eq:e36}
\end{equation}
the leading term of the metric proper-time ratio (E36); the next metric
order is smaller by another factor of $g\,\Delta h/c^2$ and is
irrelevant at any laboratory scale. The sign is pinned to the physical
statement that a higher clock runs faster: $(P-1)_{\rm grav}>0$ for
$h>h_{\rm ref}$, in the same $P-1=\Delta\nu/\nu_0$ convention every
other pivot term uses, and a shipped regression test asserts exactly
this statement. The magnitude at standard gravity is
$g/c^2=\GravPivotPerMetre$ per metre of height.

Composition follows the same additive pattern as the BBR term (E33):
$P(\mathbf{r})-1 = (P-1)_{\rm stark}(\mathbf{r}) + (P-1)_{\rm BBR} +
(P-1)_{Q}(\mathbf{r}) + (P-1)_{\rm grav}(\mathbf{r})$, with the
quadrupole term of Sec.~\ref{sec:quadrupole} present for ion species.
The gravitational term depends only on position in the potential, so
its genuine cross term with any field-driven shift is of order
$10^{-31}$, far below anything this paper resolves. One consistency
point deserves its own sentence: Eq.~\eqref{eq:e36} and the kinematic
time-dilation factor of Eq.~\eqref{eq:e21} are the two leading terms of
one proper-time expansion, the potential term and the velocity term,
and the pipeline adds each exactly once, with no $g\times v$ cross term
at any order it resolves.

The term earns its keep in the \texttt{lattice\_extended} ensemble
regime, which distributes $n_{\rm sites}$ copies of the single-site
lattice quadrature (the exact Hermite--Gauss motional average of
Sec.~\ref{sec:scalar-pivot}) along a configured axis at a configured
spacing, weighted by a Gaussian or uniform site-occupation envelope.
Every extended-lattice node is a static point exactly like a
single-site lattice node, so the fast analytic mode stays exact and the
rotor cross-check applies unchanged, and each site's own position feeds
every pivot term in scope: the local field through the Stark term, the
uniform BBR term, and Eq.~\eqref{eq:e36}'s height dependence. The
output is a per-site frequency map, the observable Bothwell et al.\
published~\cite{Bothwell2022}: each site's position, occupation weight,
and weighted mean fractional shift, together with a weighted
least-squares linear fit across the site means whose slope is the
map's headline number.

An extended sample forces one labeling discipline the single-site
regimes never needed. The frequency spread across sites is dominated by
the deterministic linear gradient: higher sites tick systematically
faster, a repeatable and in principle refocusable structure, and
folding it silently into a $T_2^*$ or linewidth would misread
deterministic inhomogeneity as stochastic decoherence. Every
extended-lattice report therefore carries a test-pinned note stating
that its ensemble spread and $T_2^*$ include the deterministic
gradient, and the site map reports both the total spread and the
gradient-removed residual spread (the weighted standard deviation of
the site means after subtracting the best-fit line), extending the
showcase's SEM-versus-$T_2^*$ discipline (Sec.~\ref{sec:showcase}) to
the deterministic-versus-stochastic axis.

Treating $g$ as uniform has a stated validity range. The error of a
single-$g$ evaluation grows quadratically with the sample's height
extent, through the height dependence of $g$ itself, and reaches the
$10^{-19}$ fractional floor only near a $\GravUniformBoundM$~m extent;
the pipeline warns at $\GravExtentWarnM$~m, an order of magnitude under
that bound, and documents that beyond laboratory scale the physically
correct input is a surveyed geopotential difference. Eq.~\eqref{eq:e36}
reports height differences within the sample relative to a local
reference, and the whole apparatus's absolute geoid and Earth-rotation
offset is a separate constant outside its scope. At the $10^{-19}$
level the choice of $g$ itself also stops being a formality: standard
gravity is a defined placeholder, and the lab's own surveyed local
value is the physically correct input. For Boulder, Colorado, the
surveyed $g=\BothwellSurveyedG$~m/s$^2$~\cite{vanWestrum2019} gives a
slope of $\BothwellLocalSlopePerMm$ per millimetre against standard
gravity's $\GravStdSlopePerMm$, a difference well inside Bothwell et
al.'s own stated prediction uncertainty (Sec.~\ref{sec:bothwell}).

\subsection{Ion-clock electric-quadrupole shift}
\label{sec:quadrupole}

Ion clocks bring a genuinely new coupling into the same pivot. A $D$-
or $F$-state upper clock level with $J\ge1$ carries a permanent
electric-quadrupole moment $\Theta$, and that moment couples to the
electric-field \emph{gradient} at first order, in contrast to the
DC-Stark and BBR terms, which are second order in the field itself.
The level shift is
\begin{equation}
  \Delta E_Q(J, m_J) \;=\; \frac{\Theta_{\rm SI}}{2}\,
  \frac{J(J+1)-3m_J^2}{J(2J-1)}\;
  \hat n^{\mathsf T} G(\mathbf{r})\,\hat n ,
  \label{eq:e34}
\end{equation}
(E34) with $\hat n$ the quantization-axis unit vector and
$G(\mathbf{r})$ the traceless symmetric part of the field-gradient
tensor $\partial_i E_j$ the field smoother already delivers. The $m_J$
factor and the overall sign follow Roos et al.'s Eq.~(1), transcribed
directly from primary text~\cite{Roos2006}, and the coordinate-free
$\hat n^{\mathsf T} G\,\hat n$ contraction is an exact rewriting of the
published axial-plus-asymmetry angular form, derived in the project's
conventions record and cross-checked numerically against that angular
form as a standing test.

For a fixed $(J, m_J)$ and quantization axis the whole angular factor
is a scalar coefficient, so the shift composes additively into the
pivot alongside every other term (E35), with no cross term: the
quadrupole coupling is first order in the gradient while the Stark and
BBR terms are second order in the field, different field quantities at
different orders. One exact identity comes free with the tensor
structure. For any three mutually orthogonal quantization axes the
three contractions sum to the trace of $G$, which is zero by
construction, so averaging a measurement over three perpendicular field
orientations cancels the quadrupole shift exactly, for any gradient
tensor; the pipeline implements this as a supported averaging mode and
tests the null at machine precision. The engine's null is the ideal
one: a real clock's residual after such averaging comes from imperfect
axis orthogonality and timing, a separate systematic the engine does
not model, and no report presents the exact cancellation as an
achievable experimental bound.

The species registry carries the measured quadrupole moments this term
needs: the $D_{5/2}$ states of Ca$^+$, Sr$^+$, and Ba$^+$ and the
$D_{3/2}$ and $F_{7/2}$ states of Yb$^+$, each with its literature
citation and verification status recorded as registry fields. The
Yb$^+$ $F_{7/2}$ entry's negative $\Theta$ doubles as the registry's
sign anchor: a shipped regression test requires it to produce an
opposite-sign shift from any positive-$\Theta$ $D$ state under an
otherwise identical gradient, $m_J$, and axis. The
$J{=}0\!\to\!J{=}0$ ion clocks Al$^+$ and In$^+$ carry no first-order
quadrupole moment at all and run through the scalar DC-Stark path of
Sec.~\ref{sec:scalar-pivot}, with a report note stating that their
small hyperfine-mediated second-order quadrupole residual is out of
scope.

Every ion-species report also carries a binding boundary statement,
because a quadrupole map alone would overstate what has been
characterized. The same stray DC field that sources the DC-Stark and
quadrupole terms also displaces an ion from its trap's RF null and
drives excess micromotion, a separate shift pathway the engine does not
model, and published measurements in a real Sr$^+$ trap
\cite{Dube2005} attribute
roughly $\IonMicromotionTensorPct\%$ of the $m_J$-dependent budget to
micromotion-driven tensor Stark shifts against roughly
$\IonQuadrupolePct\%$ from the quadrupole term alone. The report note
states this shared-cause structure explicitly, so a quadrupole map is
never presented as a stray-field budget
(Sec.~\ref{sec:limitations} places the tensor-polarizability and
micromotion package on the roadmap together for exactly this reason).

\subsection{Scope: what is and is not modeled}

\begin{table*}[htbp]
\centering
\caption{Modeled vs.\ not-modeled physics (current codebase scope).
Every "not modeled" row is a documented, intentional MVP scope boundary,
not an oversight; see Sec.~\ref{sec:limitations} and
benchmarks/MAPPING.md's engine-scope recap.}
\label{tab:scope}
\begin{tabular}{ll}
\toprule
Modeled & Not modeled \\
\midrule
DC-Stark shift, scalar $\Delta\alpha$ (Eq.~\eqref{eq:e14b}) & AC lattice-light (magic-wavelength) shift \\
Second-order (kinematic) Doppler (Eq.~\eqref{eq:e21}) & Zeeman shifts (1st or 2nd order) \\
Blackbody radiation, uniform $T$ (Eq.~\eqref{eq:e32}, $50$--$350$~K) & BBR spatial $T(\mathbf{r})$ maps, solid-angle effective temperature \\
Gravitational redshift across an extended sample (Eq.~\eqref{eq:e36}) & BBR M1/E2 multipoles ($\sim6\times10^{-20}$), hyperpolarizability \\
Ion electric-quadrupole shift, $D$/$F$-state ions (Eq.~\eqref{eq:e34}) & Tensor polarizability and RF/micromotion dynamics \\
Field-gradient inhomogeneous dephasing ($T_2^*$) & Atom-atom collisional / density shifts \\
Classical (Monte Carlo) motional ensembles & Lattice band-structure / tunneling \\
Lattice and extended-lattice (Hermite--Gauss quadrature) ensembles & Probe (clock-laser) AC Stark shift \\
$^{87}$Sr, $^{171}$Yb; Ca$^+$/Sr$^+$/Ba$^+$/Yb$^+$ $D$/$F$ states; Al$^+$/In$^+$ (scalar) & Time-dependent (RF trap / probe) fields \\
Imported (CSV/FEA) or synthetic fields & \\
\bottomrule
\end{tabular}
\end{table*}

\FloatBarrier
\section{Numerical methods}
\label{sec:methods}

\subsection{Precision discipline}
\label{sec:precision}

The target regime is fractional shifts at $10^{-17}$--$10^{-19}$,
frequently on top of stray-field biases many orders of magnitude
larger. The pipeline guards structurally
against two catastrophic-cancellation failure modes. First, the absolute
Compton phase ($\omega_C T \sim 10^{20}$~rad over a one-second
interrogation) is never accumulated; every integration path accumulates
only the \emph{perturbation} quantity $\delta\tilde\omega$ (dimensionless,
of order the fractional shift itself), with compensated (Kahan)
summation in every long accumulation loop. Second, the DC-Stark
field-squared term $|E|^2$ of Eq.~\eqref{eq:e14b} is evaluated
\emph{directly} from the total field $E(\mathbf{r})$
(\texttt{cliffordclock.integrator.omega.pivot\_perturbation\_stark}, the
rate function every scalar evaluation mode shares, with no
baseline/perturbation split supplied), not via a term-by-term expansion.
This is numerically safe for a structural reason:
unlike the absolute-Compton-phase case above, Eq.~\eqref{eq:e14b}
contains no subtraction of large near-equal quantities in the first
place. The shift is a single product chain, $|E|^2$ multiplied once by
the small prefactor $\Delta\alpha/(2h\nu_0)$, so there is no cancellation
for a term-by-term decomposition to protect against. A separate,
term-by-term decomposition of the same formula,
$|E|^2=|E_0|^2+2E_0\!\cdot\!\delta E+|\delta E|^2$ (E11's field split),
\emph{is} implemented in the library as
\texttt{stark\_pivot\_terms}, for callers that need the gradient-driven
cross term $2E_0\!\cdot\!\delta E$ isolated from the baseline as its own
quantity (e.g.\ a gradient-only systematic-budget line item); this
paper's reported shifts do not use it.

This safety claim is pinned by an adversarial
pipeline-level regression test defined in \texttt{tests/}
\texttt{test\_e2e.py} as {\small\texttt{test\_wp9\_} \texttt{major1\_}
\texttt{stark\_dc\_} \texttt{adversarial\_} \texttt{gradient\_}
\texttt{pins\_no\_} \texttt{cancellation\_} \texttt{regression}}
(recapped in \path{docs/validation.md}). The test constructs a
$|E_0|=10^5$~V/m uniform bias along with a field gradient tuned so the
shift differential between the lattice quadrature's extreme nodes is
$\sim10^{-19}$, then compares each node's directly-evaluated shift
against a 50-digit \texttt{decimal} reference: the measured per-node
absolute error is $\sim1.3\times10^{-26}$, seven orders of magnitude
below the $10^{-19}$-level signal it resolves. An embedded
discrimination guard confirms the test is not vacuous: applying the same
comparison to the genuinely dangerous antipattern (forming each node's
$O(1)$ pivot $P=1+(\text{shift})$ explicitly and recovering the
differential by subtracting two such totals) degrades to a
$\sim5\times10^{-20}$ error, which the test asserts (and confirms) must
exceed its passing bound. This also distinguishes two numbers a reader
could otherwise conflate. The $10^{-18}$ fractional-shift scale
motivating this precision discipline (Sec.~\ref{sec:intro}) is the
physical \emph{target regime}. The bounds this paper measures are
(i) relative agreement with closed-form and literature references at
$\mathrm{rtol}=10^{-10}$ or better (Table~\ref{tab:validation}) and
(ii) the absolute-error pin above, $\sim10^{-19}$--$10^{-26}$, in this
specific adversarial configuration.

These bounds span nineteen orders of magnitude and mix relative ratios
with absolute fractional-shift errors, which is exactly the condition
under which a single number gets misquoted as "the accuracy of the
tool." They are therefore never combined onto one scale in this paper,
and each carries its character wherever it is quoted: the closed-form
and literature rows of Table~\ref{tab:validation} are relative
agreements ($\mathrm{rtol}\leq10^{-10}$), the adversarial per-node
error pin is an absolute fractional-shift error
(\texttt{tests/test\_e2e.py}'s MAJOR-1 regression pin, with the same
test also recording the far larger error the
subtract-two-near-unity-pivots antipattern would produce), the
rotor/scalar agreement of Sec.~\ref{sec:rotor-verification} and
Sec.~\ref{sec:showcase} is exactly zero relative in the showcase's
ensemble mean shift, and the $10^{-18}$ physical regime is a design
target the pipeline is built for, never something this paper measured.
No one of these numbers substitutes for any other: a reader who wants
"the accuracy of a reported shift" should read the validation table
and this paper's stated target regime together, not either alone.

\subsection{Field model}

An imported or synthetic field is decomposed as
\begin{equation}
  E_{\rm total}(\mathbf{r}) = E_0(\mathbf{r}) + \delta E(\mathbf{r}) ,
  \label{eq:e11}
\end{equation}
with $E_0$ a least-squares degree-1 (uniform + linear) analytical
baseline and $\delta E$ the residual, fit by a thin-plate-spline radial
basis function per vector component (E12); the field gradient
$\partial_iE_j$ (E13) is obtained by exact automatic differentiation of
the fitted evaluator, so the field and its gradient are always
mutually consistent by construction. The RBF fit is performed with SciPy; the
gradient-consistent evaluator is re-implemented in
JAX~\cite{Jax2018} so it composes
with the pipeline's own autodiff-based downstream code
(Sec.~\ref{sec:rotor-verification}'s spin connection, and the second-order
Doppler kinematic term).

\subsection{Performance tiers within the fast analytic and trajectory modes}
\label{sec:tiers}

The scalar rate integrand $\delta\tilde\omega(\mathbf{r}(t),\mathbf{v})$
contains no Compton-frequency content (the fast $\sim10^{20}$~rad/s
carrier is removed analytically by the perturbation discipline above),
so its time variation is set entirely by the atom's own motional
timescale, microseconds to milliseconds for realistic traps. This is
what makes real (microsecond-to-second)
interrogation times tractable at all, and Table~\ref{tab:tiers} summarizes
the resulting cost structure of the three evaluation modes
(Sec.~\ref{sec:model-modes}) plus one internal performance variant within
the trajectory mode. The fast analytic mode is exact for a static
motional state in a time-independent field: every quadrature node has
$\mathbf{v}=0$ and a constant rate, so the time integral is trivial and
the cost is independent of the requested interrogation time $T$. The
trajectory mode's direct (default) variant steps the classical
integrator at a resolution set by the trap period, with
\begin{equation}
  d\tilde\tau \;=\; \frac{T_{\rm orb}}{N_{\rm res}\,\tau_C} ,
  \qquad N_{\rm res} = \StepSizeNRes~\text{(default)},
  \label{eq:e31}
\end{equation}
(E31) with $T_{\rm orb}=2\pi/\min(\omega_{xyz})$ the slowest trap axis's
period. A second trajectory-mode variant, \emph{secular averaging},
replaces per-step time integration with a single average over one
orbital period (one full cycle of the trap's own periodic motion) scaled
by the full interrogation time; it applies whenever the classical motion
is periodic in an isotropic trap. The rotor mode's
Compton-scale integrator remains available as the first-principles
cross-check every other mode/tier is validated against.

\begin{table*}[htbp]
\centering
\caption{The three evaluation modes and their performance tiers. Cost is
stated in the requested interrogation time $T$.}
\label{tab:tiers}
\footnotesize
\begin{tabular}{lllllc}
\toprule
Evaluation mode & Tier & Regime & \texttt{integration.mode} & Basis & Cost in $T$ \\
\midrule
Fast analytic & A & lattice (motional state) & fast\_path (default) & exact (E29) & $O(1)$ \\
Trajectory & B(i) & classical & direct (default) & scalar integrator, E31 step & $O(\text{steps})$ \\
Trajectory (periodic) & B(ii) & classical, isotropic trap & secular & one-orbit average (E30) & $O(1)$ \\
Rotor & C & either & worldline / Compton-fine direct & Compton-scale Cl(1,3) integrator & $O(T/\tau_C)$ \\
\bottomrule
\end{tabular}
\end{table*}

An accuracy study (Fig.~\ref{fig:step-size}, reproducing the
CONVENTIONS.md V4 closed-form harmonic-trap case: a linear-gradient
field's oscillating pivot contribution cancels exactly over full trap
periods, leaving the trap-center pivot value plus the closed-form
second-order-Doppler time average) sweeps the step size at fixed
$N_{\rm res}$ and measures a convergence order of $\StepSizeOrder$
against the exponential-midpoint stepper's design order of 2. At the
default $N_{\rm res}=\StepSizeNRes$, the phase matches the closed form
to $\StepSizeDefaultError$ relative, with rotor-norm drift
$\StepSizeDefaultDrift$, both comfortably inside this project's
$\mathrm{rtol}\le10^{-8}$ / $\mathrm{drift}<10^{-12}$ target bounds.

\begin{figure}[htbp]
\centering
\includegraphics[width=0.55\linewidth]{figures/fig3_step_size_accuracy.pdf}
\caption{This figure is a numerical convergence study of the Cl(1,3)
rotor stepper (Sec.~\ref{sec:rotor-verification}'s exponential-midpoint
integrator, E19): relative phase error against the V4
closed form, swept over the stepper's step size $d\tilde\tau$
(auto-selected via Eq.~\eqref{eq:e31}). The measured log-log slope
($\StepSizeOrder$) matches the stepper's design order of 2, and the
annotation marks that at the default step size the error is
$\StepSizeDefaultError$, comfortably below this study's own
$\mathrm{rtol}\leq10^{-8}$ tolerance. Regenerated by
\texttt{paper/figures/fig3\_step\_size\_accuracy.py} directly against
\texttt{cliffordclock.integrator.worldline}, using the same
trap/field/coupling scenario as
\texttt{notebooks/02\_step\_size\_study.ipynb} and
\texttt{tests/test\_fastpath\_select\_dtau.py}. This is a measurement of
the integrator's own convergence order, not the pipeline's real-world
shift accuracy; the bounds that actually characterize a reported
shift's precision are those of Table~\ref{tab:validation} and the
precision discipline of Sec.~\ref{sec:precision}.}
\label{fig:step-size}
\end{figure}

\subsection{Verification methodology}

Beyond the closed-form and literature comparisons of
Sec.~\ref{sec:validation}, two structural verification mechanisms
protect against a codebase that merely agrees with itself. First, the
Cl(1,3) geometric product and bivector exponential kernel are
cross-checked at the algebra level against the third-party open-source
\texttt{clifford} package~\cite{pygae-clifford} (the pygae project's
general geometric-algebra library on PyPI), used purely as an independent
test-oracle reference in CI, disjoint from this project's own
structure-tensor implementation. Second, a
plain-NumPy reference implementation
(\texttt{tests/reference\_impl.py}) re-derives the phase-accumulation
formulas directly from the equations recorded in CONVENTIONS.md,
sharing no code path, constants module, or numerical kernel with the
pipeline itself (\texttt{cliffordclock.integrator},
\texttt{cliffordclock.cl13}); the pipeline's scalar and rotor phases
are each checked against it independently (Table~\ref{tab:validation}
rows V3, E24).

\section{Validation}
\label{sec:validation}

\subsection{Labeling rules (binding)}

This section follows the labeling discipline established for this
project's benchmark work (\texttt{benchmarks/RESULTS.md},
\texttt{benchmarks/MAPPING.md}): a case is reported as
\textbf{reproducibility} when every pipeline input traces to an
independent publication and the resulting predicted value reconstructs
a number the cited source(s) already combined those same inputs to
compute; it is reported as \textbf{blind prediction} only when the
pipeline predicts a measured quantity from inputs that were not already
combined, by anyone, to produce that measured value. These are
categorically distinct claims, never conflated in this paper: a
reproducibility \textit{MET} verdict demonstrates end-to-end pipeline
correctness against independently published inputs and outputs with no
fitted or tuned parameter anywhere in the chain (a real, falsifiable
claim: a sign error, a unit-conversion bug, or a wrong
$\Delta\alpha$-to-$k_S$ derivation would show up as a predicted band
that does not overlap the published one) but it does \emph{not}
demonstrate that this engine can predict a clock shift nobody had
already computed.

Two weaker labels complete the taxonomy. An \textbf{arithmetic
reproduction} runs a published row that is itself computed, never
independently measured, back through the engine's own formula and
coefficients (Sec.~\ref{sec:jila-bbr}); a \textbf{cross-vintage
comparison} predicts a published measured quantity from an input of an
independent vintage and method, one the published result never
combined, while some other ingredient of the comparison still traces to
the source's own apparatus calibration (Sec.~\ref{sec:roos}). The
cross-vintage class sits strictly between its neighbors: weaker than a
blind prediction, because not every ingredient on both sides is
independent of every other, and stronger than an arithmetic
reproduction, because the compared-against number was never derived
from the predicting input. Neither weaker class ever joins the
reproducibility or blind-prediction counts.

As of this writing the project has two reproducibility cases
(Secs.~\ref{sec:npl} and~\ref{sec:bothwell}) and zero blind-prediction
cases; no case in this paper, this codebase's test suite, or its
benchmark reports is described as "validated against an independent
measurement" where that framing would imply the stronger, currently
unmet claim.

\subsection{Validation table}

Table~\ref{tab:validation} is generated directly by
\texttt{paper/figures/table\_validation.py} against the real pipeline,
species registry, and \texttt{benchmarks/} code. Every "Computed"
entry is a live pipeline run, and every "Reference" entry is either a
closed form evaluated independently in this script, an
independent-NumPy reference module, a literature formula applied to the
project's own cited species data, or (for the NPL row) the published
band of Sec.~\ref{sec:npl}.

\begin{table*}[htbp]
\centering
\caption{Validation summary: case, reference value/method, this
pipeline's computed value, and measured agreement. V1--V4 are
closed-form self-consistency cases (this engine against its own
analytically solvable configurations and, for V3, an independent
plain-NumPy re-implementation); KA1--KA4 are literature known-answer
cases (textbook DC-Stark and second-order-Doppler formulas, using this
project's cited $\Delta\alpha$ values); E24 is the rotor/scalar
cross-check of Sec.~\ref{sec:rotor-verification}; NPL and Bothwell are
the two experimental reproducibility cases (Secs.~\ref{sec:npl}
and~\ref{sec:bothwell}); Roos is the cross-vintage quadrupole-slope
comparison of Sec.~\ref{sec:roos}, whose \textit{NOT MET} verdict is the
expected outcome (it recovers the literature's own $\Theta$
theory-versus-measurement tension) and which joins neither headline
count. Generated by
\texttt{paper/figures/table\_validation.py}; see that script and
\texttt{tests/test\_e2e.py} / \texttt{tests/test\_known\_answers.py}
for the exact configurations reproduced.}
\label{tab:validation}
\input{generated/validation_table.tex}
\end{table*}

\subsection{The NPL reproducibility case}
\label{sec:npl}

Bowden et al.\ (NPL)~\cite{Bowden2017} independently measured a residual
stray electric field at their $^{87}$Sr lattice-clock atoms via
Rydberg-state electrometric interferometry (a completely different
atomic transition and spectroscopic method from the clock transition
itself) and converted it to a clock-transition DC-Stark shift using
the Middelmann et al.\ differential
polarizability~\cite{Middelmann2012}, the same literature value this
project's $^{87}$Sr species-registry entry cites. \textbf{No parameter
was fitted at any stage: the field magnitude and its uncertainties
(NPL's Rydberg-spectroscopy measurement), the differential
polarizability (PTB's measurement), and the clock frequency are each
independently published constants.}

NPL could have computed this number themselves, and in a sense already
did: for a spatially uniform $\NplFieldNominal$~V/m field the shift is
a single evaluation of Eq.~\eqref{eq:e14b}, so this case is not
evidence of predictive power beyond the published record. What it
demonstrates instead is that the entire pipeline (field ingestion,
lattice quadrature, precision-safe numerics, asymmetric uncertainty
propagation, and reporting) reproduces an independently published
result end-to-end with no fitted or tuned parameter anywhere in the
chain. Following this project's benchmark taxonomy
(\texttt{benchmarks/RESULTS.md}, \texttt{benchmarks/MAPPING.md}; see
the labeling rules above), that makes this a
\textbf{reproducibility} case: NPL already combined the same two
published ingredients this pipeline combines, so reconstructing their
number checks end-to-end pipeline correctness against independently
published inputs and outputs. The stronger claim, predictive power on
a measurement nobody had already converted to a shift, belongs to the
blind-prediction category, and it remains unmet.

NPL's residual field, $E=\NplFieldNominal$~V/m, carries independent,
individually asymmetric statistical and systematic uncertainties. These
are combined in quadrature separately on each side (never symmetrized
against each other, and never treated as Gaussian), giving field bounds
$[\NplFieldLo,\NplFieldHi]$~V/m. The pipeline (\texttt{coupling.type:
stark\_dc}, the lattice fast path, a genuine 1~s interrogation; the
same machinery Table~\ref{tab:validation}'s KA1 row validates) is run
three separate times at this field's low, nominal, and high bounds,
giving the predicted band
$[\NplPredictedLo,\NplPredictedHi]$, compared in
Fig.~\ref{fig:npl-band} against NPL's own published band
$[\NplPublishedLo,\NplPublishedHi]$. The two bands
\NplBandsOverlap{} (a precisely defined closed-interval overlap test);
NPL's nominal value falls inside the predicted band and vice versa.
\textbf{Verdict: \textit{\NplVerdict{}} (\NplCaseClass{} case, with no fitted
parameters anywhere in the chain).}

NPL's own single-point field measurement does not exercise this
pipeline's actual differentiator: gradient-driven shifts, inhomogeneous
broadening, and line shapes computed from a field with real spatial
structure. Eq.~\eqref{eq:e14b} at a uniform field is a one-line
calculation regardless of what pipeline performs it. Sec.~\ref{sec:showcase}
demonstrates that capability directly, end to end, on a chamber-scale
field with genuine curvature; what it does \emph{not} yet demonstrate is
that capability against a field characterization nobody has already
converted to a shift. Producing and validating that class of result
against an independently supplied field is exactly the blind-prediction
program described in Sec.~\ref{sec:limitations}, and the showcase's
scenario remains, like the NPL case above and the closed-form validation
suite, a demonstration whose every input is a documented scenario
choice.

\begin{figure}[htbp]
\centering
\includegraphics[width=0.4\linewidth]{figures/fig2_npl_band.pdf}
\caption{The NPL reproducibility case: this pipeline's predicted
DC-Stark shift band (from NPL's independently measured residual field
and this project's Middelmann-sourced $\Delta\alpha$) vs.\ NPL's own
published band~\cite{Bowden2017}. Shaded region: NPL's published band.
Error bars: asymmetric field/shift uncertainty, propagated per-side in
quadrature (no Gaussian assumption on the asymmetric field
uncertainty). Regenerated by \texttt{paper/figures/fig2\_npl\_band.py}
via \texttt{benchmarks/run\_benchmarks.run\_npl\_reproducibility\_case}.}
\label{fig:npl-band}
\end{figure}

\subsection{The Bothwell millimetre-scale redshift reproducibility case}
\label{sec:bothwell}

The second reproducibility case targets a different pivot term and a
different kind of observable. Bothwell et al.\ resolved the
gravitational redshift across a single millimetre-scale $^{87}$Sr
lattice sample, the first single-apparatus measurement of gravitational
time dilation at this length scale, publishing a per-pixel frequency
map and two independently analyzed, systematics-corrected linear
slopes~\cite{Bothwell2022}. The extended-lattice pipeline of
Sec.~\ref{sec:grav-pivot} is configured to their published sample
geometry: a Gaussian site-occupation envelope with
$\sigma=\BothwellEnvelopeSigmaUm\,\mu$m inferred from their stated
$\pm\BothwellWindowSigma\sigma$ analysis window, $\BothwellNSites$ computational sites
at the real $\BothwellSiteSpacingNm$~nm magic-lattice spacing, and the
surveyed local gravity for Boulder, $g=\BothwellSurveyedG$~m/s$^2$
(van Westrum's NOAA survey~\cite{vanWestrum2019}, the reference
Bothwell et al.\ themselves cite), the physically correct input at
this level (Sec.~\ref{sec:grav-pivot}).

Run end to end, the per-site machinery produces the map of
Fig.~\ref{fig:bothwell-sitemap} and a weighted-least-squares slope of
$\BothwellPredictedSlopePerMm$ per millimetre in Bothwell et al.'s own
coordinate convention. Their axis puts lower physical positions at
larger coordinate, so their published gradients are negative while the
engine's own height convention gives a positive slope; the benchmark
states and applies this sign mapping explicitly, so the sign agreement
is deliberate. Against the two published corrected slopes, method A's
$\BothwellMeasuredANominal\pm\BothwellMeasuredAUnc$ per millimetre (a
fourteen-dataset campaign) and method B's
$\BothwellMeasuredBNominal\pm\BothwellMeasuredBUnc$ per millimetre (a
synchronous two-region analysis), the prediction lands within
$\BothwellSigmaA\sigma$ and $\BothwellSigmaB\sigma$ respectively, and
the band verdicts are \textit{\BothwellVerdictA{}} against method A and
\textit{\BothwellVerdictB{}} against method B.

The comparison isolates exactly the term it should. The published
corrected slopes already have per-pixel density and
second-order-Zeeman corrections applied before the linear fit, plus
budget-level corrections for BBR, lattice light, DC Stark, and pixel
calibration, so the corrected number targets the isolated
gravitational gradient Eq.~\eqref{eq:e36} predicts. The benchmark's
own field is configured to exactly zero for the same reason: no Stark,
BBR, or quadrupole term is active, so the case isolates the
gravitational term to match its target. Bothwell et al.'s own DC-Stark
gradient row is in-scope physics for this engine, and it enters the
benchmark report as a narrative cross-reference only, since the
comparison target already has it corrected out.

The binding classification is \textbf{reproducibility}, with a caveat
inverted from the NPL case's. The $g/c^2$ arithmetic is textbook and
Bothwell et al.\ computed it themselves trivially, so nothing about
the coefficient demonstrates predictive power; what this case
validates end to end is the extended-sample machinery, the per-site
geometry, the envelope weighting, and the map assembly, producing the
right measured-map slope with zero adjustable inputs. NPL's caveat ran
the other way: there the differential-polarizability reconstruction
carried the domain expertise, and the pipeline's contribution was the
end-to-end evaluation. Together the two cases make the project
headline of Sec.~\ref{sec:validation}, two reproducibility cases.

\begin{figure}[htbp]
\centering
\includegraphics[width=0.7\linewidth]{figures/fig7_bothwell_sitemap.pdf}
\caption{The Bothwell reproducibility case: the engine's per-site
fractional-frequency map ($\BothwellNSites$ sites, blue points), its
weighted-least-squares slope fit (black line), and Bothwell et al.'s
two published corrected slope bands~\cite{Bothwell2022} drawn through
the same intercept (shaded wedges), all in the engine's physical-height
sign convention. The annotation restates the two sigma distances
($\BothwellSigmaA\sigma$ against method A, $\BothwellSigmaB\sigma$
against method B). Bottom panel: the Gaussian site-occupation envelope,
with the $\pm\BothwellWindowSigma\sigma$ analysis window marked. Regenerated by
\texttt{paper/figures/fig7\_bothwell\_sitemap.py} against the real
\texttt{lattice\_extended} pipeline at the exact configuration
\texttt{benchmarks/run\_bothwell\_redshift.py} pins.}
\label{fig:bothwell-sitemap}
\end{figure}

\subsection{A weaker case: the Aeppli et al.\ (JILA) BBR arithmetic-reproduction check}
\label{sec:jila-bbr}

The BBR term (Sec.~\ref{sec:bbr-pivot}) carries a deliberately weaker
check than Sec.~\ref{sec:npl}'s NPL case: whether the engine
arithmetically reproduces a BBR shift a source group already computed
from its own measured temperature and coefficients. Aeppli et al.\
report their $^{87}$Sr clock's in-vacuum resistance-thermometer (RTD)
temperature, $T=\JilaBbrTemperature\pm\JilaBbrTemperatureUnc$~K, and
separately publish the BBR-row shift it produces~\cite{Aeppli2024}.
Running Eq.~\eqref{eq:e32} at that temperature with the registry
coefficients gives $(P-1)_{\rm BBR}=\JilaBbrPredictedNominal$, with a
combined coefficient-plus-temperature band of
$[\JilaBbrPredictedLo,\JilaBbrPredictedHi]$. Aeppli et al.'s published value
spans $[\JilaBbrPublishedLo,\JilaBbrPublishedHi]$, leaving a residual of
$\JilaBbrResidual$; the bands \JilaBbrBandsOverlap{} (verdict
\textit{\JilaBbrVerdict}).

\textbf{This agreement is expected almost by construction}: the
registry's Sr-87 dynamic coefficients are anchored to Aeppli et al.'s
own value through the same PTB rescaling
Sec.~\ref{sec:bbr-pivot} describes, so the case reproduces a paper its
own inputs trace to. Nothing about Aeppli et al.'s shift was independently
measured the way NPL's field was. Their published row is itself an
arithmetic evaluation from their own inputs. Following this project's binding
labeling discipline (\texttt{benchmarks/RESULTS.md},
the project's theory sign-off record (G7), B5), that makes this an
\textbf{arithmetic reproduction of a published standard-formula
evaluation}: weaker than this paper's reproducibility cases, and weaker
still than the (unmet) blind-prediction category. It does \emph{not}
join the validation headline, now two reproducibility cases, NPL and
Bothwell (Secs.~\ref{sec:npl} and~\ref{sec:bothwell}). What the case does establish is that the
engine's BBR arithmetic and coefficient-provenance chain are correct:
any $10^{-19}$-class number here is \textbf{arithmetic-reproduction
fidelity, not BBR accuracy}.

\subsection{A cross-vintage comparison: the Roos two-ion quadrupole slope}
\label{sec:roos}

The quadrupole term of Sec.~\ref{sec:quadrupole} carries its own
benchmark case, against the measurement that anchors the field. Roos
et al.\ prepared an entangled two-ion Ca$^+$ state, measured its
parity-oscillation frequency as a function of an applied, mechanically
calibrated axial field gradient, and fit a line with slope
$a=\RoosMeasuredSlopeCompact$~Hz\,mm$^2$/V
\cite{Roos2006}. The engine predicts that slope from Itano's
independent theory value
$\Theta=\RoosThetaTheory\,ea_0^2$~\cite{Itano2006} through its own
factor chain: the $m_J$ factors of Eq.~\eqref{eq:e34} evaluated by
real engine calls, the doubling of the gradient each ion sees from the
presence of the other, and the resulting two-ion enhancement over a
single ion's shift, whose engine-derived ratio of
$\RoosEnhancementRatio$ structurally reproduces Roos et al.'s stated
$24/5$.

The predicted slope is $\RoosPredictedSlope$~Hz\,mm$^2$/V, a residual
of $\RoosResidualPct\%$ against the measured value, and the band test
returns \textit{\RoosCrossVintageVerdict}. That verdict is the expected
outcome, and the case record says so in a self-contextualizing note it
carries in its own JSON: the residual reproduces the literature's own
tension between Itano's computed
$\Theta=\RoosThetaTheory\,ea_0^2$ and Roos et al.'s measured
$\Theta=\RoosThetaMeasuredCompact\,ea_0^2$, so the
comparison validates the engine's factor chain while isolating the
coefficient the two vintages disagree on.

A second, deliberately circular variant closes the loop on the
arithmetic. Feeding the engine Roos et al.'s own extracted
$\Theta=\RoosThetaMeasuredCompact\,ea_0^2$, the very
value their paper derives by inverting the same Fig.~4a fit through
$\Theta=(5/12)\,h\,a$, predicts
$\RoosArithPredictedSlope$~Hz\,mm$^2$/V, a $\RoosArithResidualPct\%$
residual whose bands \RoosArithBandsOverlap{} (\textit{\RoosArithVerdict}). This
variant is labeled an arithmetic reproduction, circular by
construction, and its value is the factor-consistency check it
performs: the engine's independently derived coefficient chain
reproduces Roos et al.'s own published conversion relation exactly,
confirmed by a closed-loop round trip at float64 precision.

The headline variant carries the taxonomy's cross-vintage-comparison
label (Sec.~\ref{sec:validation}'s labeling rules), and the label
keeps its two weaknesses visible: Roos et al.'s applied gradient is
itself calibrated through their own trap model, and their own fit
produced the very slope being predicted. It is nonetheless a genuine
external comparison, distinct from the circular variant above, because
Itano's $\Theta$ was never an ingredient of the fit it predicts.
Neither variant joins the paper's headline counts, and the project
remains at two reproducibility cases and zero blind predictions.

\subsection{The covered slice of one platform's published budget}
\label{sec:budget-slice}

The validation cases above were built one systematic at a time; this
subsection reads them the way a lab would, as rows of one platform's
evaluation. The JILA Sr system publishes both of the evaluations this
paper reproduces from published inputs, so Table~\ref{tab:budget-slice}
collects the covered rows of that single platform side by side with
the lab's own numbers, each row in one unit convention, each published
value carrying its stated uncertainty, and each comparison reduced to
a difference and a distance in units of that uncertainty. Two readings
carry the table. The BBR difference sits at $\BudgetBbrSigma\sigma$ of
the published uncertainty ($\BudgetBbrSigmaCombined\sigma$ against the
published and prediction bands combined in quadrature), small by
construction since the registry's
dynamic coefficients are anchored to this group's own measurement,
which is exactly why that row keeps the weaker arithmetic-reproduction
label. The two redshift slopes are Bothwell et al.'s own independent
analysis methods, and their central values differ from each other by
$\BudgetAbSpread\times10^{-20}$/mm ($\BudgetAbSigma\sigma$ combined),
more than either differs from this pipeline's single prediction, which
lies between them. The one in-scope row without a comparison, their
DC-Stark gradient, is exactly the shape of the blind-prediction
program of Sec.~\ref{sec:limitations}: comparing against it requires a
residual-field characterization the paper does not publish, and a lab
sharing that characterization turns this context row into the record's
first blind prediction. The composed, executable form of this table is
\path{notebooks/11_real_budget_slice.ipynb}, which also prints the
report's own exclusion notes naming the rows this release does not
model.

\begin{table*}[htbp]
\centering
\caption{The covered slice of one platform's published evaluation (the
JILA Sr system: Aeppli et al.~\cite{Aeppli2024}, Bothwell et
al.~\cite{Bothwell2022}), computed from published inputs by the same
case objects the benchmark suite commits. $\sigma$ is the difference in
units of the published value's stated uncertainty, a conservative
convention that ignores this tool's own prediction band; for the BBR
row, the only one where that band is comparable to the published one,
the combined-uncertainty metric reduces the distance to
$\BudgetBbrSigmaCombined\sigma$. The context row is
in-scope physics (Eq.~\eqref{eq:e14b}) awaiting the lab-side field
characterization described in Sec.~\ref{sec:limitations}. Regenerated
by \texttt{paper/figures/table\_budget\_slice.py}; the executable form
is \texttt{notebooks/11\_real\_budget\_slice.ipynb}.}
\label{tab:budget-slice}
\footnotesize
\input{generated/budget_slice_table.tex}
\end{table*}

\FloatBarrier
\section{Showcase: a chamber-scale field's inhomogeneous shift budget}
\label{sec:showcase}

This section is the pitch the rest of the paper builds toward: from a
CAD/FEA-style field export with genuine spatial structure to the full
inhomogeneous-shift budget of a real atomic cloud (the mean shift
together with the spread across the ensemble, the resulting dephasing
time $T_2^*$, and the clock transition's line profile), computed by the same
pipeline validated in Sec.~\ref{sec:validation}, with the trajectory and
rotor evaluation modes checked directly against each other on the
identical ensemble. This is a demonstration scenario, built the same way
Sec.~\ref{sec:npl}'s NPL case and this paper's closed-form validation
suite are: every input is a chosen, documented scenario parameter, not
sourced from an independent field measurement (Sec.~\ref{sec:limitations}
states what that stronger case would still require). The scale is
anchored to a documented class of real events: Lodewyck et al.\
observed and then cancelled a stray-charge DC-Stark shift at an
operating Sr lattice clock~\cite{Lodewyck2012}, and a field event of
this scenario's magnitude would dominate the clock's coherence
outright, with nothing in the unmodeled physics competing at that
scale. The scenario therefore answers the two questions a practitioner
actually brings to it: what signature a charge-patch event would leave
on an operating clock's shift, spread, and $T_2^*$, and what
stray-field level a chamber geometry under design can tolerate before
it matters.

\subsection{Scenario: an asymmetric two-electrode chamber with a patch spot}

The field is a finite-difference solution of Laplace's equation
(\texttt{examples/generate\_showcase\_field.py}; the same from-scratch,
NumPy/SciPy-only conjugate-gradient solver as this project's earlier
COMSOL-format round-trip test, with no COMSOL or other paid software
involved) for a grounded chamber containing two electrode plates of
\emph{different} size and voltage at different, non-mirror-symmetric
positions, plus a small circular patch-potential spot embedded in one
wall (a crude but standard model of trapped-charge contamination on a
chamber surface). Unlike a simple parallel-plate capacitor (whose field is
nearly linear across any small interior region), this scenario is chosen
specifically to carry real curvature and cross-terms: the exported field
region shows a $\ShowcaseRelFieldVariation\%$ relative magnitude
variation over one cloud standard deviation (0.82~mm, two field-grid
cells at this scenario's 0.41~mm spacing) from the trap center along its
steepest axis. This is a genuinely varying field. Figure~\ref{fig:showcase}(a)
shows the whole-chamber field magnitude with both electrodes and the
patch spot marked; the field at the trap center itself is
$|E|=\ShowcaseFieldCenter$~V/m.

\textbf{Scale reasoning (why the cloud is chamber-scale, not
lattice-site scale).} A real optical-lattice-site cloud's motional
ground-state extent is tens of nanometers (the harmonic-oscillator
scale $\sqrt{\hbar/(2m\omega)}$, at a typical lattice-site trap
frequency $\omega\sim10^5$~rad/s), far smaller than any chamber-scale
field feature, so it would never sample genuine field curvature; a
field-gradient map's real spatial content only shows up in
an ensemble whose own spatial extent is a meaningful fraction of the
field's variation scale. This showcase therefore uses a classical,
chamber-scale trapping stage (e.g.\ a large-volume "science" dipole or
magnetic trap) with a $\ShowcaseTrapOmega$~rad/s
isotropic trap and a $\ShowcaseTemperatureUK\,\mu$K, genuinely
laser-cooled-scale $^{87}$Sr ensemble ($\ShowcaseEnsembleSize$ classical
Monte Carlo atoms), giving a classical thermal cloud size
$\sigma=\sqrt{k_BT/(m\omega^2)}=\ShowcaseCloudSigmaUm\,\mu$m, a sizable
fraction of the field's own spatial-variation scale, so the ensemble's
trajectories genuinely sample real field curvature. This is a scenario-\emph{geometry} choice
(trap frequency and temperature), documented exactly as such; no physical
constant is touched to reach it.

\subsection{Trajectory and rotor modes, run head to head}

Each atom's Sr-87 classical trajectory (Maxwell--Boltzmann velocities,
velocity-Verlet propagation in the trap, over a genuine
$\ShowcaseInterrogationTime$~s window spanning several trap oscillation
periods) is propagated through the fitted field and accumulates its
DC-Stark phase via the pipeline's \textbf{trajectory mode}
(Sec.~\ref{sec:model-modes}), exactly as
\texttt{examples/showcase\_gradient\_\allowbreak dispersion\_sr87.yaml}
runs through the standard \texttt{cliffordclock} CLI/API. The identical Monte Carlo
trajectories are then re-accumulated through the \textbf{rotor mode}'s
Cl(1,3) engine, driven directly against the same trajectories (the
shipped config schema's \texttt{integration.mode=worldline} keyword
exposes the rotor cross-check for static lattice nodes only,
Sec.~\ref{sec:model-modes}; the identical accumulator function
also runs, unchanged, against a moving classical trajectory, which is
what this figure exercises). \textbf{The two evaluations agree to
$\ShowcaseMeanShiftRelDiff$ relative in the ensemble mean shift, to
$\ShowcaseMaxPhaseDiff$ absolute (float64-noise scale) in every
individual atom's accumulated phase, and to $\ShowcaseMaxShiftDiff$
absolute in every individual atom's fractional shift}. That is the
message this section exists to make concrete: the trajectory mode is not an
approximation of a "more correct" rotor calculation, and the rotor mode
is not required to
get a spatially varying field's shift budget right; both are already
the same number, verified on this exact scenario as well as on the
simpler cases of Table~\ref{tab:validation}.

\subsection{The inhomogeneous shift budget}

The pipeline reports a mean fractional shift
$\langle\Delta\nu/\nu_0\rangle=\ShowcaseMeanShift$ with a population
spread of $\ShowcaseShiftSpread$ across the ensemble (SEM
$\ShowcaseShiftSem$). The spread is the physically meaningful
dispersion quantity here; the SEM in parentheses is only the
statistical precision of the mean. The ensemble's accumulated phase spread $\sigma_\Phi$ implies
an inhomogeneous dephasing time via the standard Gaussian-dephasing form
\begin{equation}
  T_2^* = \frac{\sqrt{2}\,T}{\sigma_\Phi} ,
  \label{eq:e27}
\end{equation}
(E27, CONVENTIONS.md), giving
$T_2^*=\ShowcaseTtwoStarUs\,\mu$s, roughly three orders of magnitude
below this scenario's own interrogation window. An experiment run this
way would lose coherence to field-gradient inhomogeneity long before
the interrogation time itself, a real, quantitative finding a lab could
not get from a field-gradient map and a mean-shift calculation alone.
(This number is set by where each atom sits in the field: the mean
phase and its spread both scale with the interrogation window here, so
$T_2^*$ reflects the ensemble's spatial distribution through the field
regardless of the particular window chosen; see
\texttt{examples/showcase\_gradient\_dispersion\_sr87.yaml}'s
own scale-reasoning note.) Figure~\ref{fig:showcase}(c) shows the
resulting per-atom shift distribution and spectral line profile, with
$T_2^*$ annotated; the rotor-mode line profile is overlaid and
indistinguishable from the trajectory-mode one at this resolution
(rotor $T_2^*=\ShowcaseTtwoStarRotorUs\,\mu$s).

\textbf{What a lab could not otherwise easily get.} A general-purpose
FEA solver reports a field map; converting that map into a
gradient-driven dispersion budget for a specific trapped ensemble
(spread, $T_2^*$, and line shape, beyond a mean shift at one point) is
exactly the workflow gap Sec.~\ref{sec:intro} describes, and exactly what
this section demonstrates end to end, with the scalar and rotor engines
agreeing on the answer.

\begin{figure*}[htbp]
\centering
\includegraphics[width=0.95\linewidth]{figures/fig4_showcase_gradient_dispersion.pdf}
\caption{The showcase scenario (Sec.~\ref{sec:showcase}). (a) Whole-chamber
field magnitude at the trap's $z$-plane, with the two asymmetric
electrodes and the patch-potential wall spot marked (dashed outlines,
projected onto this plane; see the electrode/patch $z$-positions
annotated on the panel). (b) The locally fitted field (the same
\texttt{FieldSmoother} fit the pipeline evaluates) with a subsample of
Monte Carlo trajectories overlaid, each colored by its final accumulated
fractional shift. (c) The per-atom shift distribution (histogram) and the
resulting spectral line profile (both trajectory- and rotor-mode
evaluations overlaid), with $T_2^*$ annotated. Regenerated by
\texttt{paper/figures/fig4\_showcase\_gradient\_dispersion.py} against the
real pipeline, using \texttt{examples/showcase\_gradient\_dispersion\_sr87.yaml}
and its generating field \texttt{examples/showcase\_field.txt}
(\texttt{examples/generate\_showcase\_field.py}).}
\label{fig:showcase}
\end{figure*}

\FloatBarrier
\section{Limitations and roadmap}
\label{sec:limitations}

The current physics package covers the field-driven and relativistic
terms of an optical-clock frequency budget end to end: DC-Stark shifts
from static and gradient fields, electric-quadrupole shifts for ion
clocks, blackbody radiation at a uniform temperature, second-order
Doppler, and the gravitational redshift across an extended atomic
sample. Every one of these terms runs through the same numerics
Sec.~\ref{sec:methods} pins, so a lab can add a new field geometry or
species and get a full dispersion budget from the same validated
machinery. Table~\ref{tab:scope} states the remaining boundary
directly, and every item outside it is a documented, intentional
choice with a stated reason. An independent benchmarking pass against
two published clock uncertainty budgets~\cite{Aeppli2024,Jia2026}
found most published systematic line items outside this project's
scope, with blackbody radiation dominating both budgets by three or
more orders of magnitude; that finding is what ordered the scope
extensions shipped so far, and it orders the queue below.

Lattice light shifts are the next physics package on that queue,
because they complete the lattice-clock systematic triad alongside BBR
and the field-driven terms already covered. Modeling them at the magic
wavelength needs species-specific AC polarizability data: the
hyperpolarizability, the E2/M1 polarizabilities, and the
motional-state intensity-sampling dependence that recent dual-ensemble
evaluations resolve directly. Tensor polarizability and RF/micromotion
dynamics for ion traps arrive together as one package, because the
published Sr$^+$ measurements of Sec.~\ref{sec:quadrupole}
\cite{Dube2005} show
micromotion-driven tensor Stark shifts dominating the
$m_J$-dependent budget; shipping tensor polarizability without the
time-dependent RF field that drives micromotion would model the
smaller term while implying full coverage of the larger one. That same
time-dependent field machinery, built once, also sets up probe-beam AC
Stark support broadly.

Two further packages wait on partner-lab demand. Density and
collisional shifts are a different physics class altogether: they need
an interaction model between atoms, with s-wave and p-wave collision
physics and excitation-fraction dependence, where this project's
ensembles are independent particles by construction; some labs manage
this systematic entirely through density control and extrapolation, in
which case it may stay outside this tool's scope indefinitely.
Blackbody radiation's spatial refinement, a $T(\mathbf{r})$ map
reduced to a solid-angle effective temperature, is a natural next step
once a partner lab wants to hand over thermal FEA output directly
instead of pre-reducing it to the single number the uniform-$T$ term
takes today.

The rotor engine implements the physical DC-Stark coupling directly (the
$\Omega$-bivector construction \texttt{build\_omega\_stark},
Sec.~\ref{sec:rotor-verification}), with its own dedicated, direct verification
against the scalar path (including a re-run of the NPL reproducibility
case through the rotor integrator and the showcase's head-to-head Monte
Carlo comparison, Sec.~\ref{sec:showcase}). By design, the fast analytic
and trajectory modes' performance tiers (\texttt{fast\_path},
\texttt{direct}, \texttt{secular}; Table~\ref{tab:tiers}) stay on the
scalar accumulator: they are built around a coupling-agnostic rate
function, and routing them through a per-step Cl(1,3) rotor construction
would trade away their performance advantage for no accuracy benefit,
given the scalar/rotor equivalence Sec.~\ref{sec:model-modes} establishes
for every physical coupling this project currently implements. That
equivalence also locates the rotor engine's roadmap value precisely: it
begins where a future extension's physics stops being expressible as
the single scalar pivot this section's other
modes assume. Magnetic couplings (Zeeman shifts, Table~\ref{tab:scope}),
internal-state structure beyond a two-level clock transition, and
transport-induced geometric phases (a rotor accumulated along the
moving trajectory itself) are the
concrete candidates, and each needs the rotor's bivector/spinor
structure to be represented at all.

The clearest concrete next step toward a genuine \textbf{blind
prediction} case (Sec.~\ref{sec:validation}) is access to a published
source that independently characterizes a stray field this project's
own pipeline has not already been used, by that same source, to convert
to a shift. One such candidate is known by citation
(Li et al.\ 2024~\cite{Li2024}, cited by Jia et al.\ 2026~\cite{Jia2026}
as the source that originally characterized an applied field underlying
one of their systematic-shift budget entries) but was not accessible to
this project through any legitimate open-access route at the time of
writing; its content remains unknown and is neither ruled in nor ruled
out as a future reproducibility or blind-prediction case. A structured
collaboration with a clock-apparatus group willing to supply a field
characterization ahead of (or independent of) their own shift
measurement, a beta blind-prediction program, is the concrete path
to closing this gap, and is future work; no such collaboration is
named here pending the owner's own outreach.

\FloatBarrier
\section{Code availability and citation}
\label{sec:code}

The source code, test suite, documentation, and the figure-generation
scripts underlying this paper are available at
\url{https://github.com/velar-mbr/CliffordClock} under the GNU
Affero General Public License v3.0 or later (AGPLv3). The Python
package and command-line interface are both distributed under the name
\texttt{cliffordclock} (pip package and CLI command). A machine-readable
citation file (\texttt{CITATION.cff}) is provided in the repository
root.

\bibliography{refs}

\end{document}
